\chapter{Deep Learning}

\section*{What the Last Three Chapters Were Waiting On}

Chapters 5 through 7 each ran into some version of the same wall. Naive Bayes worked from hand-counted words. Support vector machines worked from a hand-chosen kernel. Bayesian networks worked from a hand-drawn graph and hand-specified probability tables. In every case, a human decided, before training began, what the relevant features or structure of the problem were, and the learning algorithm only ever operated on that human decision, never on the raw material underneath it.

Chapter 4 already showed a way around this, in miniature, without anyone noticing at the time that it was the way around this. A hidden layer, trained with backpropagation, doesn't just fit a boundary. It builds its own internal representation, one the original input space didn't hand it, and it does this without a person specifying what that representation should look like. Nothing in principle stops you from stacking more hidden layers on top of each other, letting each one build a representation out of the layer before it: a first layer might learn something simple, a second layer might learn something built out of combinations of the first layer's output, and so on, each layer a little more abstract than the one beneath it. If that worked reliably, it would resolve the feature-engineering bottleneck outright. Nobody would need to decide in advance which words, which kernel, which graph structure mattered. The network would work that out, layer by layer, directly from raw data.

This is deep learning's core promise, and the idea itself is not new to this chapter. Backpropagation, from Chapter 4, already tells you how to compute the gradients needed to train a network with any number of layers. The genuinely strange part of this chapter is that the idea sat mostly unused for nearly two decades after backpropagation became widely known in the mid-1980s. Deep networks, more than two or three hidden layers, were tried, and they mostly didn't train well. The reason turns out to be a specific, measurable problem, not a vague sense that "deep networks are hard."

\section*{Watching Gradients Vanish}

Build a network with ten hidden layers, all the same width, using the sigmoid activation function from Chapter 4, and run a single forward pass followed by a single backward pass, exactly the computation Chapter 4's training loop repeats thousands of times. Record the size of the weight gradient at every layer along the way, from the layer closest to the output back to the layer closest to the input.

\begin{Verbatim}[fontsize=\small, frame=single, xleftmargin=1em, xrightmargin=1em]
def forward_backward(layers, x, activation, dactivation):
    activations = [x]
    a = x
    for W, b in layers:
        z = a @ W + b
        a = activation(z)
        activations.append(a)

    delta = activations[-1] - 0.0
    grad_norms = []

    for i in reversed(range(len(layers))):
        W, b = layers[i]
        a_out = activations[i+1]
        a_in = activations[i]
        delta = delta * dactivation(a_out)
        grad_W = a_in.reshape(-1, 1) @ delta.reshape(1, -1)
        grad_norms.append(np.linalg.norm(grad_W))
        delta = delta @ W.T

    grad_norms.reverse()
    return grad_norms
\end{Verbatim}

With identical random initialization, differing only in whether the activation function is sigmoid or ReLU (\texttt{relu(x) = max(0, x)}, a function this chapter hasn't used yet but is about to matter a great deal):

\begin{Verbatim}[fontsize=\small, frame=single, xleftmargin=1em, xrightmargin=1em]
A 10-hidden-layer network, 20 units wide, identical random
initialization, differing only in activation function.

Gradient norm at each layer, layer 0 = closest to the input:

 layer        sigmoid           relu
     0     0.00129359   40410.079108
     1     0.00180970   36283.608235
     2     0.00352348   20390.571518
     3     0.00923468   12729.186098
     4     0.02438025   13018.023637
     5     0.05288814    9725.821551
     6     0.11537179    5613.132937
     7     0.20535379    8724.612728
     8     0.47538652    6617.408564
     9     1.09062402    7210.266542

Ratio of layer-0 gradient to layer-9 gradient (closest-to-input vs closest-to-output):
  sigmoid: 0.0011861032
  relu:    5.6045194543
\end{Verbatim}

Look at the sigmoid column first. The gradient at layer 9, the layer nearest the output, is a reasonable 1.09. By layer 0, the layer nearest the input, ten layers of backpropagation later, it has shrunk to 0.0013, less than a thousandth of its starting size. The ratio at the bottom makes it exact: the input-side gradient is roughly one nine-hundredth the size of the output-side gradient. That's not a training run that happens to be going slowly. It's a training run in which the earliest layers of the network are receiving a learning signal too small to meaningfully update their weights, no matter how long you wait. Add more layers and the effect compounds; each additional layer multiplies in another factor smaller than one, because the sigmoid function's derivative is never larger than 0.25 anywhere and is close to zero across most of its range. Stack enough layers and the gradient reaching the earliest ones underflows into a value indistinguishable from noise.

Now look at the ReLU column. The numbers are much larger overall, because ReLU doesn't squash its input into a narrow range the way sigmoid does, but more importantly, the ratio between layer 0 and layer 9 doesn't collapse toward zero. If anything it grows slightly, but it stays within the same rough order of magnitude across all ten layers rather than shrinking by three orders of magnitude. The gradient signal reaches the earliest layers largely intact.

\section*{Why This Explains a Twenty-Year Gap}

This single property, whether a network's activation function preserves or destroys gradient magnitude across many layers, goes a long way toward explaining why deep networks were mostly considered impractical from the mid-1980s until roughly the mid-2000s to early 2010s, despite backpropagation itself having been available and well understood that entire time. Researchers weren't failing to try deep networks. They were training them with sigmoid and tanh activations, the standard choice inherited from Chapter 4, and running directly into the vanishing gradient problem this section just measured directly. The earliest layers, the ones responsible for learning the simplest, most foundational features everything else would be built on top of, were the layers receiving the weakest training signal, which is close to the worst possible arrangement for a system that's supposed to learn hierarchically.

The eventual fixes arrived in stages, not all at once. Geoffrey Hinton and colleagues showed in 2006 that training a deep network one layer at a time, in an unsupervised pretraining phase, before fine-tuning the whole stack together, could get around the worst of the vanishing gradient problem well enough to train networks meaningfully deeper than had been practical before. This was an important proof that depth could work at all, though it was a workaround rather than a resolution: it sidestepped the vanishing gradient rather than removing its cause. The more direct fix came from simply changing the activation function. Rectified linear units, ReLU, output exactly their input when positive and exactly zero when negative, and unlike sigmoid, their derivative is either exactly 1 or exactly 0, never something small and shrinking. Combined with better-designed random weight initialization schemes tuned specifically to keep signal variance roughly constant from layer to layer, ReLU largely dissolved the vanishing gradient problem for networks of the depth researchers actually wanted to train, as the comparison above shows directly.

\section*{What This Actually Unlocks}

It's worth being precise about what got resolved here and what didn't. This chapter has not yet shown a network learning features from raw data the way the opening section promised; that demonstration, edge detectors and shape detectors emerging on their own from raw pixels, belongs to the next chapter, where the specific architecture built for that purpose, the convolutional network, finally makes it concrete. What this chapter has shown is narrower and more foundational: the specific mechanical obstacle that made "just add more layers" impractical for two decades, made visible and measured directly, and the specific, comparatively small change, a different activation function plus more careful initialization, that removed it.

That narrowness is worth sitting with. The representational bottleneck from Chapter 2 was resolved by an architectural idea, adding a hidden layer. The feature-engineering bottleneck from Chapters 5 through 7 is resolved, in principle, by the same idea taken further: stack many hidden layers and let each build on the last. But between having the right idea and having it actually work sat a purely practical, computational obstacle that had nothing to do with whether the idea was correct. Backpropagation was correct in 1986. Depth was the right direction in 1986. Neither fact made deep networks trainable until the activation functions and initialization schemes caught up, decades later. That specific gap, between a correct idea and a trainable one, is this chapter's claim, and it's a mechanical, historical fact about depth in particular: what changed between 1986 and 2012 was optimization catching up to representation, nothing broader. Whether that same shape of gap, a correct idea outrunning the practical means to realize it, turns out to be a recurring feature of this whole field rather than a one-time story about depth, is a question worth holding until the very last chapter, once there's a full history's worth of evidence to check it against.
